\begin{abstract}
In interactive imitation learning, a learner practises a task and an expert takes
over to demonstrate the right actions. Every takeover costs expert time, which
compute cannot reduce. Under a fixed budget, what each demonstration contains
decides the result. Existing methods decide only when to take over. They call the
expert at the first step of an episode where the policy becomes uncertain. Two
decisions are left to chance: which failed episode to correct, and where the
demonstration begins. A failed episode runs from start to end without achieving
the task. We present \textbf{DISEIL} (\textbf{D}emonstration d\textbf{I}stillation
for \textbf{S}ample-\textbf{E}fficient \textbf{I}mitation \textbf{L}earning), which
decides both. After each round of practice, DISEIL groups the failed episodes into
modes of failure. For each mode it prescribes where a new demonstration should
begin, and checks the start is reachable before any expert time is spent. We
evaluate on five tasks, with image and state observations, ten settings in all,
at a budget of 20 demonstrations, against the best baselines. DISEIL reaches the
highest mean success rate in all ten, on average about 3 percentage points above
the strongest baseline.
\end{abstract}
